\section{Conclusion}
\label{sec:conclusion}

This report has documented the complete arc of a memory-augmented incident-triage project: from dataset construction (instrumenting ten microservices under a strict no-label-leakage discipline, running 100 fault-injection experiments, building a 347-ticket humanized Jira memory corpus and 110-ticket distractor pool), through four method generations (V1 numeric baselines, V2 multi-channel skill chain, V2-advanced LLM-and-graph pipelines, and the final Tiered Cascade Hybrid), to eight late-stage refinement experiments and a meta-analysis identifying the project's actual headline.

The technical core is a four-layer cascade. L1 stacks six pipeline triage scores via class-balanced logistic regression. L2 fuses four retrievers via Reciprocal Rank Fusion, anchored by a multi-retriever overlap-rerank for top-1 precision. L3 detects novel failure modes through three OR-combined signals: an LLM agent verification stage, a retrieval-confidence-threshold free signal, and a learned per-window logistic regression added in late-stage refinement. L4 composes the final output. The cascade runs at $16~\mu$s per window at inference and on 1{,}008 in-distribution test windows achieves Hit@5 $= 0.912$, Hit@1 $= 0.722$, MRR $= 0.794$, strict PR-AUC $= 0.9998$, and novel-recall $= 0.793$ at novel-precision $= 0.940$ (a $+388\%$ relative improvement over the locked baseline at preserved precision). It is robust to memory noise (Hit@5 drops only $2\%$ at $50\%$ distractor ratio) and to family-distribution shift (novelty F1 within $11\%$ of in-distribution under leave-one-family-out cross-validation).

\textbf{Three findings transcend any single experiment.} (1) The cascade is composition-fragile, not component-fragile: a $+404\%$ rel standalone retrieval lift broke cascade integration through score-scale assumptions that are currently implicit. Future cascade extensions must address score normalization explicitly. (2) Three independent rerankers (cross-encoder, agent-rerank, LLM judge) failed in sequence; multi-retriever consensus on top-1 dominates single-LLM judgment on this dataset. (3) The retrieval channel was already at $93.5\%$ of its theoretical union ceiling; the real headline win came from an underexploited novelty channel, not from retrieval refinement. We discovered the win rather than designed it---a finding about the original cascade's blind spots rather than a clever architecture.

\textbf{The contribution to the field is twofold.} A reproducible, instrumented dataset paired with the construction pipeline that produced it---bridging a real gap between public observability datasets and Jira corpora. And a concrete worked example of a four-layer cascade that composes seven Pareto-incomparable pipelines, with documented evidence of where the composition holds and where it breaks. We hope both will be useful for further research on memory-augmented incident operations and for engineering teams considering similar systems in production.

\paragraph{Acknowledgments.} Online Boutique is courtesy of Google's microservices-demo project. The TAWOS reference dataset was used for length/style anchoring only. Qwen 3.6 35B-A3B was served locally via LM Studio with MoE offload. All compute was performed on a single workstation with an RTX 5060 (8 GB VRAM).

\paragraph{Availability.} All scripts, scenario YAMLs, fault-injection harness, derived feature definitions, pipeline implementations, the cascade builder, the G-series experiments, and the per-phase observation documents are in the project repository. Reproduction instructions are in \texttt{docs3/19-REPRODUCE.md}; the locked dataset is at \texttt{data/derived/global/2026-05-25-dataset-v5-large-global/}; per-phase observations under \texttt{docs4/}.
